% Standalone problem document, generated from the adjacent README.md.
% Compile with XeLaTeX. Mathematical content is maintained in Markdown.
\documentclass[11pt,a4paper]{article}
\usepackage[fontsize=10.5pt]{scrextend}
\usepackage[margin=22mm,headheight=15pt,headsep=9mm,footskip=11mm]{geometry}
\usepackage{fontspec}
\setmainfont{texgyrepagella}[Extension=.otf,UprightFont=*-regular,BoldFont=*-bold,ItalicFont=*-italic,BoldItalicFont=*-bolditalic]
\setsansfont{texgyreheros}[Extension=.otf,UprightFont=*-regular,BoldFont=*-bold,ItalicFont=*-italic,BoldItalicFont=*-bolditalic]
\setmonofont{texgyrecursor}[Extension=.otf,UprightFont=*-regular,BoldFont=*-bold,ItalicFont=*-italic,BoldItalicFont=*-bolditalic,Scale=MatchLowercase]
\usepackage{amsmath,amssymb}
\usepackage{unicode-math}
\setmathfont{texgyrepagella-math.otf}
\usepackage{xcolor,microtype,enumitem,needspace,fancyhdr}
\usepackage{bookmark}
\definecolor{ink}{HTML}{17344B}
\definecolor{accent}{HTML}{197D80}
\definecolor{muted}{HTML}{596773}
\hypersetup{colorlinks=true,linkcolor=accent,urlcolor=accent,pdftitle={TR-19 - Exact
best-rank-one approximation ratios for general tensor
formats},pdfauthor={Open Problems in Numerical Linear Algebra}}
\urlstyle{same}
\setlength{\parindent}{0pt}
\setlength{\parskip}{5pt plus 1pt minus 1pt}
\linespread{1.04}
\setlength{\emergencystretch}{3em}
\widowpenalty=10000
\clubpenalty=10000
\displaywidowpenalty=10000
\allowdisplaybreaks[1]
\setlist{nosep,leftmargin=1.5em}
\providecommand{\tightlist}{\setlength{\itemsep}{0pt}\setlength{\parskip}{0pt}}
\providecommand{\pandocbounded}[1]{#1}
\pagestyle{fancy}
\fancyhf{}
\fancyhead[L]{\sffamily\footnotesize\color{muted}OPEN PROBLEMS IN NUMERICAL LINEAR ALGEBRA}
\fancyhead[R]{\sffamily\bfseries\footnotesize\color{accent}TR-19}
\fancyfoot[L]{\parbox[b]{0.9\textwidth}{\sffamily\fontsize{8}{10}\selectfont\color{muted}Editorial ratings; a bounded search cannot certify that a problem remains unsolved.\\Literature check: 2026-09-08}}
\fancyfoot[R]{\sffamily\footnotesize\color{muted}\thepage}
\renewcommand{\headrulewidth}{0.3pt}
\renewcommand{\headrule}{\hbox to\headwidth{\color{accent}\leaders\hrule height \headrulewidth\hfill}}
\makeatletter
\renewcommand{\section}{\Needspace{5\baselineskip}\@startsection{section}{1}{0pt}{12pt plus 2pt minus 2pt}{4pt}{\sffamily\bfseries\large\color{ink}}}
\renewcommand{\subsection}{\Needspace{4\baselineskip}\@startsection{subsection}{2}{0pt}{9pt plus 1pt minus 1pt}{3pt}{\sffamily\bfseries\normalsize\color{ink}}}
\makeatother
\setcounter{secnumdepth}{0}
\begin{document}
{\sffamily\small\color{muted}Tensor computations\par}
\vspace{3pt}
{\raggedright\hyphenpenalty=10000\exhyphenpenalty=10000\sffamily\bfseries\fontsize{21}{24}\selectfont\color{ink}Exact
best-rank-one approximation ratios for general tensor formats\par}
\vspace{7pt}
{\sffamily\small\textbf{Difficulty:} extreme\quad\textbf{Importance:} interesting
to the community\par}
\vspace{4pt}
{\color{accent}\hrule height 0.8pt}
\vspace{6pt}
\textbf{Topic:} Tensor approximation and norm comparison.

For integers \(d\ge3\) and \(2\le n_1\le\cdots\le n_d\), let
\(V=\mathbb R^{n_1\times\cdots\times n_d}\). For \(T\in V\), define

\[
\|T\|_F^2=\sum_{i_1=1}^{n_1}\cdots\sum_{i_d=1}^{n_d}T_{i_1\ldots i_d}^2,
\]

\[
\|T\|_\sigma=\max_{\substack{x_j\in\mathbb R^{n_j}\\\|x_j\|_2=1,\ 1\le j\le d}}
\left|\sum_{i_1=1}^{n_1}\cdots\sum_{i_d=1}^{n_d}
T_{i_1\ldots i_d}\prod_{j=1}^d(x_j)_{i_j}\right|.
\]

Determine, as a function of the dimensions, the exact value

\[
\tau(n_1,\ldots,n_d)=\min_{T\in V\setminus\{0\}}\frac{\|T\|_\sigma}{\|T\|_F}.
\]

Equivalently, determine the largest constant \(\tau\) for which
\(\tau\|T\|_F\le\|T\|_\sigma\) holds for every tensor in that format.
The target is the sharp constant for general formats, not merely another
upper or lower bound. All formats constitute one problem.

This ratio measures the worst-case amount of tensor energy accessible to
a best rank-one approximation and is relevant to convergence analyses of
successive rank-one approximation methods.

\subsection{References}\label{references}

\begin{enumerate}
\def\labelenumi{\arabic{enumi}.}
\tightlist
\item
  Z. Li and Y.-B. Zhao,
  \href{https://pure.port.ac.uk/ws/portalfiles/portal/19117904/On_norm_compression_inequalities.pdf}{On
  norm compression inequalities for partitioned block tensors},
  \emph{Calcolo} 57 (2020), article 11, §5, equation (15), pp.~15--16.
  The text explicitly identifies determination of the general extremal
  ratio as the outstanding problem.
\item
  Z. Li, Y. Nakatsukasa, T. Soma, and A. Uschmajew,
  \href{https://arxiv.org/abs/1707.02569}{On orthogonal tensors and best
  rank-one approximation ratio}, \emph{SIAM Journal on Matrix Analysis
  and Applications} 39 (2018), 400--425, abstract and characterization
  of equality in the basic dimension bound.
\item
  K. Kozhasov and J. Tonelli-Cueto,
  \href{https://arxiv.org/abs/2201.02191}{Probabilistic bounds on best
  rank-one approximation ratio}, arXiv v2 (2022-09-23), abstract and
  bounds for general and partially symmetric tensors.
\end{enumerate}

\subsection{Status check --- 2026-09-08}\label{status-check-2026-09-08}

Checked the 2020 published text and the later probabilistic-bound
paper's latest listed arXiv version. Searches for ``best rank-one
approximation ratio 2026'' and ``extremal ratio tensor spectral
Frobenius 2025 2026'' found no exact general formula. Known
orthogonal-tensor cases and probabilistic dimension bounds give partial
answers. They are not additional open entries. This target differs from
counting critical rank-one approximations, choosing an inner product to
minimize that count, and tensor-train quasi-optimality. No recent
explicit reaffirmation of the full exact-value problem was located.
\end{document}
